\documentclass[../../main.tex]{subfiles}
% \documentclass[a4paper,10pt]{ltjsarticle}
\begin{document}
\rhead{京都大学後期$1996$年$6$番}
\setcounter{equation}{0}
\begin{oframed}
  \documentclass[../../main.tex]{subfiles}
\begin{document}

$n$を$3$以上の整数とする．
円周上の$n$等分点のある点を出発点とし，$n$等分点を一定の方向に次のようにすすむ．
各点でコインを投げ，表が出れば次の点に進み，裏が出れば次の点を飛び越しその次の点にすすむ．

\begin{enumerate}
  \item 最初に$1$周回ったとき，出発点を飛び越す確率$p_n$を求めよ．
  \item $k$は$2$以上の整数とする．
    $k-1$周目までは出発点を飛び越し，$k$周目に初めて出発点を踏む確率を$q_{n,k}$とする．
    このとき$\lim_{n\to\infty}q_{n,k}$を求めよ．
\end{enumerate}

\end{document}

\end{oframed}
\setlength{\columnseprule}{0.4pt}
\begin{multicols}{2}
  {\bf [解答]}

  \paragraph{(1)}

  \begin{figure}[H]
    \centering
    \begin{tikzpicture}[scale=1.25,>=stealth,every node/.style={font=\small}]
      % n=7 の例．番号は進行方向に増える．
      \draw[very thick,blue!85!black] (0,0) circle (1.65);
      \foreach \k in {0,...,6}{
        \coordinate (S\k) at ({90-360*\k/7}:1.65);
        \node[circle,fill=blue!85!black,inner sep=1.3pt] at (S\k) {};
        \node at ({90-360*\k/7}:2.0) {$\k$};
      }
      \draw[->,very thick,teal!85!black] (S0) to[bend right=15] (S1);
      \draw[->,very thick,red!85!black] (S0) to[bend right=15] (S2);
      \draw[->,very thick,orange!90!black] (S6) to[bend left=15] (S1);
      \node[above=7pt] at (S0) {出発点};
      \node[teal!85!black,right] at (1.3,1.3) {表：$+1$};
      \node[red!85!black,right] at (1.3,0.95) {裏：$+2$};
    \end{tikzpicture}
    \caption{$n=7$の例．$6\to1$では出発点$0$を飛び越す}
    \label{1996-6:fig:1}
  \end{figure}

  $p_{n+2}$ について$1$回目に出たコインで場合分けすると
  \begin{itemize}
    \item $1$回目に表で，出発点をとびこすとき，$p_{n+1}$
    \item $1$回目に裏で，出発点をとびこすとき，$p_n$
  \end{itemize}
  だから，
  \begin{align}
    p_{n+2} = \dfrac{p_{n+1} + p_n}{2} \label{eq:1}
  \end{align}
  である．初期条件は
  \begin{align}
    p_3 & =\dfrac{3}{8}  \\
    p_4 & =\dfrac{5}{16}
  \end{align}
  である．以下\cref{eq:1}を解く．変形して
  \begin{align}
    p_{n+2}+\dfrac{1}{2}p_{n+1} & = p_{n+1}+\dfrac{1}{2}p_{n}
  \end{align}
  だから初期条件の元で解いて，$n\ge 3$に対して
  \begin{align}
    p_{n+1} = -\dfrac{1}{2}p_{n}+\dfrac{1}{2}
  \end{align}
  となる．さらに変形すると
  \begin{align}
    p_{n+1}-\dfrac{1}{3} = -\dfrac{1}{2}\left(p_{n}-\dfrac{1}{3}\right)
  \end{align}
  だから，$n\ge 3$に対して
  \begin{align}
    p_n
    & = \dfrac{1}{3} + \dfrac{1}{6}\left(\dfrac{-1}{2}\right)^{n-1} \\
    & = \dfrac{1}{3}\left(1-\left(\dfrac{-1}{2}\right)^{n}\right)
  \end{align}
  が求める答えである．$\cdots$(答)

  \vspace{10pt}
  \paragraph{(2)}
  $k-1$周目までとびこす時，まず$1$回目をとびこす確率は$p_n$である．
  $2$周目以降でとびこす確率は，とびこした後の点が$1$周目の出発点より$1$つだけ進んだ点であることから$p_{n-1}$である．
  又，$k$回目に出発点を踏む確率は排反を考えて$(1-p_{n-1})$である．
  以上から，$\alpha=-\dfrac{1}{2}$とおいて
  \begin{align}
    q_{n,k}
    & =  p_n \cdot (p_n)^{k-2} \cdot (1-p_{n-1}) \\
    & =  \dfrac{1}{3}(1-\alpha^n) \left\{ \dfrac{1}{3}(1-\alpha^{n-1}) \right\}^{k-2} \cdot \left(\dfrac{2}{3}+\dfrac{1}{3}\alpha^{n-1}\right) \\
    & =  \left(\dfrac{1}{3}\right)^{k}(1-\alpha^n)(1-\alpha^{n-1})^{k-2} \left(2+\alpha^{n-1}\right)                 \\
    &\xrightarrow{n\to\infty}  2\left(\dfrac{1}{3}\right)^k
  \end{align}
  が求める極限である．$\cdots$(答)

  \vspace{10pt}
  {\bf [解説]}
  \newpage
\end{multicols}
\end{document}
